\section{Data and Methods}
\label{app:data}

\paragraph{Reproducibility.} Every number in the text, tables and figures is
generated by the scripts in the accompanying repository from public sources,
cached on first download with a SHA-256 checksum. Running \texttt{make all}
rebuilds the paper; \texttt{data/SOURCES.md} lists each dataset's URL, access
date and licence.

\paragraph{Country groups.} Small Island Developing States (SIDS) and
Landlocked Developing Countries (LLDC) follow the UN M49 classification
\parencite{unm49}. \enquote{Pacific} denotes SIDS in Melanesia, Micronesia and
Polynesia. Of \nSIDSunList{} SIDS in the M49 list, \nSIDS{} appear in World Bank
data and are used here. Group statistics are unweighted medians across
countries, so that no single large member dominates.

\paragraph{Remoteness.} For country $i$,
\begin{equation}
  R_i \;=\; \sum_{j\neq i} \frac{Y_j}{\sum_{k\neq i} Y_k}\, d_{ij},
\end{equation}
where $Y_j$ is GDP in current US dollars (\windowStart--\windowEnd{}
average, World Bank) and $d_{ij}$ is the great-circle distance between the
most populous cities of $i$ and $j$ from CEPII's GeoDist \parencite{mayer2011}.
GeoDist (2011) uses some superseded country codes; successor states inherit
their predecessor's distances (e.g. Montenegro from Serbia and Montenegro;
South Sudan from Sudan; Curaçao and Sint Maarten from the Netherlands
Antilles). Taiwan has distances but no World Bank GDP, so it carries no
weight. Remoteness is computed for \nCountriesRemote{} economies.

\paragraph{Regression.} \Cref{tab:regression} reports OLS with
heteroskedasticity-robust (HC1) standard errors. All variables are
\windowStart--\windowEnd{} averages; government effectiveness is the
Worldwide Governance Indicators estimate. Coefficients on log$_{10}$ variables
are elasticities; the \enquote{doubling} effect is $2^{\beta}-1$.

\paragraph{Shipping.} UNCTAD's quarterly LSCI is averaged to calendar years
(complete years only). \enquote{Hubs} are the ten highest-scoring economies in
\lsciFirstYear{}, fixed at the start to avoid selecting on later outcomes.

\paragraph{Cables.} From TeleGeography's public map data, a cable counts
for a country if it is in service (not planned, ready-for-service year up to
2025) and lands in at least two countries. Branch systems are recorded by
TeleGeography under their trunk's ready-for-service year.

\paragraph{Technology adoption.} For each CHAT technology, adoption is the
first year with a positive value. The frontier year is the year the fifth
country adopted; the lag is the median across later adopters. The intensity
gap is the ratio of median per-capita use in the richest quarter of
countries (CHAT real GDP per capita) to that in the poorest half, in the last
year up to 2000 in which at least 40 countries report. Technologies whose
records begin simultaneously for many countries (left-censored) are excluded.
This is a descriptive analogue of \textcite{comin2010,comin2018}, not a
replication.

\paragraph{Other variables.} Physicians per 1,000: latest World Bank value
2013--2023. Diabetes: NCD-RisC age-standardised prevalence, adults 18+,
simple mean of male and female estimates. Graduate emigration: IAB, 2010,
population 25+, 20 OECD destinations. Displacement: IDMC new disaster
displacements via the World Bank, 2008--2023, per 1,000 residents, with
unreported years counted as zero. Emissions: CO$_2$ excluding land use, tonnes
per person (EDGAR via the World Bank).

\paragraph{Departures from the original plan.} The ND-GAIN vulnerability
index and ITU broadband prices could not be downloaded automatically when the
pipeline was built (their servers refused scripted access), and the WHO
Global Health Observatory API was unavailable; WorldRiskIndex was considered
but rejected because its exposure component scales with absolute population
and so ranks large countries as most exposed. The climate figure therefore
uses displacement per head, and diabetes uses NCD-RisC. See \texttt{TODO.md}.
